\subsection{三角関数と弧度法}
$\sin x$ を定義から微分しようとすると,加法定理から
\[
 \frac{\sin(x+h)-\sin x}{h}
 =\sin x\cdot\frac{\cos h-1}{h}
   +\cos x\cdot\frac{\sin h}{h}
\]
となる.したがって,\,$\frac{\sin h}{h}$ と $\frac{\cos h-1}{h}$ の極限が分かればよい.まずは実験をしてみよう.角を弧度法で測り,\,$h$ を小さくして $\frac{\sin h}{h}$ を計算すると,次のようになる.
\begin{center}
\begin{tabular}{ccccc}
\toprule
$h$ & $1$ & $0.5$ & $0.1$ & $0.01$\\
\midrule
$\frac{\sin h}{h}$ & $0.8415$ & $0.9589$ & $0.9983$ & $0.99998$\\
\bottomrule
\end{tabular}
\end{center}
$1$ に近づきそうである.一方,角を度数法で測り,\,$h$ 度の正弦を $h$ で割ると,\,$h=1,0.1,0.01$ のいずれでも約 $0.01745$ となり,\,$1$ にはならない.この値の意味は後で分かる.以後,三角関数を微分するときは,角は必ず弧度法で表す.

予想した極限
\[
 \lim_{h\to0}\frac{\sin h}{h}=1
\]
を証明しよう.\,$0<h<\frac\pi2$ として,単位円の中心を $O$,点 $(1,0)$ を $A$,角 $h$ に対応する円上の点を $P$ とする.半直線 $OP$ と直線 $x=1$ の交点を $Q$ とすれば,\,$P=(\cos h,\sin h)$,\,$Q=(1,\tan h)$ である.図\ref{図:正弦の極限}のように,三角形 $OAP$ は扇形 $OAP$ に含まれ,扇形は三角形 $OAQ$ に含まれる.面積を順に書くと
\[
 \frac12\sin h<\frac h2<\frac12\tan h
\]
だから,\,$\sin h<h<\tan h$ を得る.

\begin{figure}[htbp]
\centering
\begin{tikzpicture}[x=3cm,y=3cm,>=Stealth]
 \fill[softblue] (0,0)--(1,0)--(1,.83910)--cycle;
 \fill[softteal] (0,0)--(1,0) arc[start angle=0,end angle=40,radius=1]--cycle;
 \draw[->] (-.05,0)--(1.3,0) node[right]{$x$};
 \draw[thick] (1,0) arc[start angle=0,end angle=90,radius=1];
 \draw[thick] (0,0)--(1,.83910)--(1,0);
 \draw[thick] (1,0)--(.766044,.642788);
 \draw (0.25,0) arc[start angle=0,end angle=40,radius=.25];
 \node at (.34,.12) {$h$};
 \fill (.766044,.642788) circle(.6pt) node[above left]{$P$};
 \node[above right] at (1,.83910) {$Q$};
 \node[below] at (1,0) {$A$};
 \node[below left] at (0,0) {$O$};
\end{tikzpicture}
\caption{三角形,扇形,外側の三角形の面積を比べる.}
\label{図:正弦の極限}
\end{figure}

$\sin h<h$ を $h$ で割り,\,$h<\tan h$ の両辺に $\frac{\cos h}{h}$ を掛けると,
\[
 \cos h<\frac{\sin h}{h}<1
\]
となる.\,$h\to0+$ で $\cos h\to1$ だから,挟み撃ちの原理から右極限は $1$ である.\,$\frac{\sin h}{h}$ は $h$ の符号を変えても同じ値だから,左極限も $1$ である.

もう一方の極限は,半角の公式 $1-\cos h=2\sin^2\frac h2$ を使って,今示した極限に帰着させてやる.
\begin{align*}
 \frac{\cos h-1}{h}
 &=-\cfrac{2\sin^2\cfrac h2}{h}\\
 &=-\sin\frac h2\cdot\cfrac{\sin\cfrac h2}{\cfrac h2}\\
 &\longrightarrow0\cdot1=0
\end{align*}
これらを最初の式に代入すれば,
\[
 \frac{\sin(x+h)-\sin x}{h}\longrightarrow\cos x
\]
を得る.同様に余弦を微分でき,商の公式から正接も求まる.
\[
 (\sin x)'=\cos x,\qquad
 (\cos x)'=-\sin x,\qquad
 (\tan x)'=\frac1{\cos^2x}\quad(\cos x\ne0).
\]

最後に,度数法の実験で現れた値を説明しておこう.\,$x$ 度は $\frac{\pi x}{180}$ ラジアンなので,\,$x$ 度の正弦は $\sin\frac{\pi x}{180}$ である.合成関数の微分から
\[
 \frac{\dd}{\dd x}\sin\frac{\pi x}{180}=\frac{\pi}{180}\cos\frac{\pi x}{180}
\]
となり,換算係数 $\frac{\pi}{180}\approx0.017453$ が現れる.実験の約 $0.01745$ はこの値だったのだ.上の簡潔な公式が成り立つのは,弧度法を使っているからである.
